The square-function argument needs a uniform $L^p$ bound for the sequence of smooth annular cutoffs supplied as inputs to the dilated operators.

\begin{Proposition}[Annular Littlewood--Paley bound]\label{ext:stein-annular-square}
Let $n\ge1$, $1<p<\infty$, and let $q\in C_c^\infty(\mathbb R^n)$ be supported in $\{c\le|\xi|\le C\}$ with $0<c\le C<\infty$. Use
\[
 \widehat f(\xi)=\int_{\mathbb R^n}e^{-ix\cdot\xi}f(x)\,dx,\qquad
 \check q(x)=(2\pi)^{-n}\int_{\mathbb R^n}e^{ix\cdot\xi}q(\xi)\,d\xi .
\]
For $j\in\mathbb Z$, define $P_j=q(2^{-j}D)$ on Schwartz functions by
$\widehat{P_j f}(\xi)=q(2^{-j}\xi)\widehat f(\xi)$. On $L^p(\mathbb R^n)$ use the convolution formula $P_j f=2^{jn}\check q(2^j\cdot)*f$. Then
\[
 \left\|\left(\sum_{k\in F}|P_{k+\ell} f|^2\right)^{1/2}\right\|_{L^p(\mathbb R^n)}
 \le C_{n,p,q}\|f\|_{L^p(\mathbb R^n)}
\]
for every finite $F\subset\mathbb Z$, every integer shift $\ell$, and every $f\in L^p(\mathbb R^n)$. The bound also holds for the full sum. The undilated kernel $\check q$ is Schwartz and has integral zero because $q(0)=0$. Its decay and that of its first derivatives give the kernel bound in \cite{r-32}, Theorem~6.1.2, (6.1.3), so (6.1.4), p.~421, applies. Restriction to the indices $F+\ell$ preserves the bound.
\end{Proposition}
